\begin{pastexam}[label={px:ntut-109-ci-04}]{北科大資工所 109 學年度 第 3 題}
（簡答）

Please consider the following questions and answer them concisely. (6 pts)
\begin{enumerate}[label=(\arabic*)]
\item Solve the \textit{continuous-knapsack problem} for the following weights ($w$), profits ($p$),
and knapsack capacity ($M$). Please show your work step by step. (3 pts)
\begin{center}
$w_1=25$, $w_2=15$, $w_3=20$, $w_4=10$, $w_5=24$;\\
$p_1=10$, $p_2=3$, $p_3=5$, $p_4=8$, $p_5=8$;\\
$M=75$.
\end{center}
\item Consider two tennis players, A and B, playing a series of games until one of them win $n$
games. Assume that the probability of player A winning a game is the same for each game and equals
to $p$, and let the probability of player A losing a game is $q=1-p$. (Hence, there are no ties.)
Let $P(i,j)$ be the probability of player A winning the series if player A needs $i$ more games to
win the series and player B needs $j$ more games to win the series. Set up a recurrence relation for
$P(i,j)$ that can be used by a dynamic programming algorithm to find the probability of player A
winning a seven-game series. (3 pts)
\end{enumerate}
\end{pastexam}

\begin{pastsoln}
答案：(1) $x = (x_1, \ldots, x_5) = (1, 0, \tfrac{4}{5}, 1, 1)$，總利潤 $30$。
(2) $P(i, j) = p\,P(i-1, j) + q\,P(i, j-1)$（$i, j \ge 1$），$P(0, j) = 1$（$j \ge 1$），
$P(i, 0) = 0$（$i \ge 1$）；七戰四勝要的是 $P(4, 4)$。

\textbf{(1)} 分數背包可以只放一部分，照資結書第~19.2~節的貪心法：依單位重量的利潤 $p_i / w_i$
由大到小放，放得下就整個放，最後一個放不下的放一部分（資結書定理~19.2 證明這樣最佳）。 先算比值：
\begin{center}
\begin{tabular}{cccccc}
\toprule
物品 $i$ & 1 & 2 & 3 & 4 & 5 \\
\midrule
$w_i$ & 25 & 15 & 20 & 10 & 24 \\
$p_i$ & 10 & 3 & 5 & 8 & 8 \\
$p_i / w_i$ & $0.4$ & $0.2$ & $0.25$ & $0.8$ & $0.333\ldots$ \\
\bottomrule
\end{tabular}
\end{center}
順序是 $4, 1, 5, 3, 2$。 逐步放：
\begin{center}
\begin{tabular}{clccc}
\toprule
步驟 & 動作 & 放的比例 & 剩下容量 & 累計利潤 \\
\midrule
1 & 物品 4（$w = 10$） & $1$ & $75 - 10 = 65$ & $8$ \\
2 & 物品 1（$w = 25$） & $1$ & $65 - 25 = 40$ & $18$ \\
3 & 物品 5（$w = 24$） & $1$ & $40 - 24 = 16$ & $26$ \\
4 & 物品 3（$w = 20$），只剩 $16$ & $16/20 = 4/5$ & $0$ & $26 + \tfrac{4}{5} \cdot 5 = 30$ \\
5 & 物品 2 & $0$ & $0$ & $30$ \\
\bottomrule
\end{tabular}
\end{center}
解是 $x_1 = 1$、$x_2 = 0$、$x_3 = 4/5$、$x_4 = 1$、$x_5 = 1$，總重量
$25 + 16 + 10 + 24 = 75$ 剛好裝滿，最大利潤 $30$。

\textbf{(2)} 看下一局。 A 還差 $i$ 勝、B 還差 $j$ 勝（$i, j \ge 1$）時，下一局：
\begin{itemize}
\item 機率 $p$ A 贏，變成 A 還差 $i - 1$ 勝、B 還差 $j$ 勝；
\item 機率 $q$ A 輸，變成 A 還差 $i$ 勝、B 還差 $j - 1$ 勝。
\end{itemize}
每局的結果互相獨立，之後 A 贏得系列賽的機率只跟「還差幾勝」有關，所以由全機率
\[
P(i, j) = p\,P(i - 1, j) + q\,P(i, j - 1), \qquad i \ge 1,\ j \ge 1 .
\]
邊界：A 已經不差勝場（$i = 0$）而 B 還差（$j \ge 1$）時，A 已經贏了，$P(0, j) = 1$；反過來
$P(i, 0) = 0$（$i \ge 1$）。 $P(0, 0)$ 不會出現，因為兩人不會同時拿到第 $n$ 勝。

七戰四勝制是 $n = 4$，所求是 $P(4, 4)$。 DP 時開一個 $(n + 1) \times (n + 1)$ 的表，先填邊界，再依
$i$ 由小到大、同一個 $i$ 裡 $j$ 由小到大填（每一格只用到左邊和上面那一格，資結書第~18.1.3~節的
由下而上）：

\begin{pseudo}{系列賽 A 獲勝的機率}
\Function{Series}{$n, p$}
    \State $q \gets 1 - p$
    \For{$j \gets 1$ \To $n$}
        \State $P[0, j] \gets 1$
    \EndFor
    \For{$i \gets 1$ \To $n$}
        \State $P[i, 0] \gets 0$
        \For{$j \gets 1$ \To $n$}
            \State $P[i, j] \gets p \cdot P[i-1, j] + q \cdot P[i, j-1]$
        \EndFor
    \EndFor
    \State \Return $P[n, n]$
\EndFunction
\end{pseudo}

共 $n^2$ 格、每格 $O(1)$，時間 $\Theta(n^2)$、空間 $\Theta(n^2)$。 舉例：$p = 1/2$ 時由對稱得
$P(4, 4) = 1/2$；$p = 0.6$ 時 $P(4, 4) = 11097/15625 \approx 0.710$，每局只多一成的勝率，七戰四勝的
勝率就到七成。 直接寫成遞迴不做記憶化的話，同一個 $(i, j)$ 會被算很多次，呼叫次數是指數級的，
這正是需要 DP 的原因。
\end{pastsoln}
